\subsection{Первообразная функции на интервале. Теорема о структуре множества первообразных. Неопределённый интеграл и его свойства}

\subsubsection*{1. Первообразная функция}

\begin{definition}
Функция $F(x)$ называется \textbf{первообразной} функции $f(x)$ на интервале $(a, b)$, если для всех $x \in (a, b)$ выполняется:
\[ F'(x) = f(x). \]
\end{definition}

\textbf{Пример.} $F(x) = \sin x$ является первообразной функции $f(x) = \cos x$ на $\R$.

\subsubsection*{2. Теорема о структуре множества первообразных}

\begin{theorem}
Если $F(x)$ — некоторая первообразная функции $f(x)$ на интервале $(a, b)$, то любая другая первообразная этой функции имеет вид $F(x) + C$, где $C \in \R$ — произвольная константа.

Иными словами, множество всех первообразных функции $f(x)$ на $(a, b)$ есть однопараметрическое семейство функций $\{F(x) + C \mid C \in \R\}$.
\end{theorem}

\begin{consequence}
Две первообразные одной и той же функции на интервале отличаются на константу.
\end{consequence}

\subsubsection*{3. Неопределённый интеграл}

\begin{definition}
\textbf{Неопределённым интегралом} функции $f(x)$ называется семейство всех её первообразных:
\[ \int f(x)\,dx = F(x) + C, \quad C \in \R, \]
где $F(x)$ — произвольная фиксированная первообразная функции $f(x)$.
\end{definition}

Символ $\int$ называется \textbf{знаком интеграла}, $f(x)$ — \textbf{подынтегральной функцией}, $f(x)\,dx$ — \textbf{подынтегральным выражением}.

\subsubsection*{4. Основные свойства неопределённого интеграла}

\begin{enumerate}
    \item \textbf{Дифференцирование интеграла:}
    \[ \left(\int f(x)\,dx\right)' = f(x), \qquad d\!\int f(x)\,dx = f(x)\,dx. \]

    \item \textbf{Интегрирование дифференциала:}
    \[ \int F'(x)\,dx = F(x) + C, \qquad \int dF(x) = F(x) + C. \]

    \item \textbf{Однородность (вынесение константы):} для любого $\alpha \neq 0$
    \[ \int \alpha\, f(x)\,dx = \alpha \int f(x)\,dx. \]

    \item \textbf{Аддитивность:}
    \[ \int \bigl(f(x) \pm g(x)\bigr)\,dx = \int f(x)\,dx \pm \int g(x)\,dx. \]

    \item \textbf{Инвариантность формулы интегрирования} относительно замены переменной: если
    $\int f(x)\,dx = F(x) + C$, то для любой дифференцируемой функции $x = \varphi(t)$
    \[ \int f(\varphi(t))\,\varphi'(t)\,dt = F(\varphi(t)) + C. \]
\end{enumerate}

\subsubsection*{5. Связь с производной и дифференциалом}

Неопределённый интеграл и операция дифференцирования взаимно обратны друг другу (с точностью до константы):
\[ \frac{d}{dx}\int f(x)\,dx = f(x), \qquad \int df(x) = f(x) + C. \]
